% New English translation from checked cote 25, site 109 (PDF page 110).
% AI-translated by GPT-6 Astra (OpenAI, ChatGPT Pro mode).
\scanpage{25/109}
\mspage{12}

in view of par.25 on the \unsure{resultant|resultant} and the discriminant of forms, will be the
deter\var{\over{m}{m}}{m}{writing over}{ega25_25_109-v0001}ination of relative divisors $\hand{\geqslant} 0$ on a multiprojective bundle.
\var{\inserted[typed]{2.8.}}{2.8.}{insertion}{ega25_25_109-v0002}
If $\underline{D}$ is a relative divisor on P, we define its multidegree as being
% NOTE: the right edge of the source page cuts off the final letters of "étan[t]" and "ramené[e]".
that of $\underline{O}_P(D)$. As above, the determination of $\mathrm{Div}^{+}(P/S)$ reduces
to that of $\mathrm{Div}^{\underline{n}}(P/S)$, for a given multidegree $\underline{n}\hand{\in}\underline{\unsure{N|M}}^{I}$, which provides an
isomorphism
\[ \underline{L} =\underline{O}_X(D) \hand{\simeq} \underline{M} \otimes_{\underline{O}_S} \underline{O}_P(\underline{n}) \ . \]
But by Chap II par.2 we obtain (\var{\xout{or by Künneth}}{}{typewriter strike-out}{ega25_25_109-v0003}
\[ (xx) \qquad f_{\hand{*}}(\underline{D}) = \underline{M} \otimes f_{\hand{*}}(\underline{O}_P(\underline{n})) = \underline{M} \otimes \hand{\bigotimes}_i \underline{\mathrm{Sym}}^{n_i}(\underline{E}_i) \quad . \]
Proceeding now as for 2.2. we find

\uline{Proposition} 2.9. With the preceding notation, to be recalled, we have\var{\xout{\illegible{1}}}{}{typewriter strike-out}{ega25_25_109-v0004} a
can one-to-one corr between the set $\mathrm{Div}^{\underline{n}}(P/S)$ of relative divisors of
multidegree $\underline{n}$ on P, and the set of invertible quotient Modules $\underline{M}$
of $\hand{\bigotimes}_i \underline{\mathrm{Sym}}^{n_i}(\underline{E}_i)^{\hand{\vee}}$ (or equivalently ...). If D and $\underline{M}$ correspond,
$\underline{O}_P(D)$ is can isomorphic to $\underline{M}\otimes\underline{O}_P(\underline{n})$, and $s_D$ is then identified
with what one guesses, taking account of (xx).

\uline{Corollary} 2.10. The subfunctor $\mathrm{Div}^{\underline{n}}_{P/S}$ of $\mathrm{Div}^{+}_{P/S}$ corresponding to
relative divisors of multidegree $\underline{n}$ is can representable by the projective
bundle $P(\hand{\bigotimes}_i \underline{\mathrm{Sym}}^{n_i}(\underline{E}_i)^{\hand{\vee}})$, and $\underline{\mathrm{Div}}^{+}_{P/S}$ is can repr by the prescheme
sum of these \unsure{latter|latter}, for $\underline{n}\hand{\in}\underline{N}^{I}$.

2.11. The preceding very simple determination of the $\mathrm{Div}_{P/S}$ is due to the
very simple structure of Pic(P/S) (in fact, to the ``discrete'' structure of the
Picard prescheme $\underline{\mathrm{Pic}}_{P/S}$ ....). We can abstract the reasoning
carried out, which amounts essentially to establishing a ``relative'' representability
(with respect to $\underline{\mathrm{Pic}}$). F$_o$r this, let us be given a morphism
\[ f\colon X \hand{\to} S \]
proper and flat of finite presentation, \uline{and} an invertible Module $\underline{L}$ on X.
\tnote{25-109-1}{The source endings ant or aht in résultant, and rni or mmi in derniers, give the same respective English words; the doubts remain marked. Corr biun can abbreviates canonical one-to-one correspondence, while can and repr abbreviate canonically and representable. The subscript o in Pour is retained in For. The X subscript in the first formula and D as the argument of the direct image in (xx) are preserved. The parenthesis before the canceled Künneth reference remains open.}

\mspage{109}
